% Pursuing Stacks (cotes 134-2 to 134-8): archive cote 134-2, site page 24 (PDF page 25).
% letter to Quillen. Draft compared: ps1.tex:470-536.
% AI-transcribed by Claude Opus 5.5 (claude-opus-5-5, adaptive thinking, effort high), Message Batches API, batch ps1 (msgbatch_01Rv5hCaNTtUt4JHx1yqoHPL), 2026-11-12; not yet checked by eye.
\scanpage{134-2/24}
\mspage{6} \margin{10}

obvious sense by a finite number,  nor even by a finite number of infinit\unsure{e}
series of maps such as $k^{i}_{l}$ , $inv_{i}$ , $\overset{i}{\underset{\hand{l}}{*}}$ , the compatibility arrows in
the pentagon,  and the like.  More precisely still,  I am convinced that
none of the functors $\underline{C}_{n} \to \underline{C}_{n+1}$ is an equivalence,  which amounts to say-%
ing that  the structures of increasing order form a\xout{n} strictly increasing
sequence  -  at every step,  there is actual extra structure added.  This
is perhaps evident beforehand to topologists in the know,  but I confess
that for the time being it isn't to me,  in terms uniquely of the some-%
what formal description I will make of the passage of $\underline{C}_{n}$ to $\underline{C}_{n+1}$ \hand{(${}^{10}$)}.
\deleted{which} \inserted{This step} theoretically is all that remains to be done,  in order to achieve \margin{\unsure{NB that $\overset{i}{\underset{\ell}{*}}$ is of order $\ell$ is heuristically clear, but will require a proof none the less}}
an explicit construction of the \xout{\unsure{notion of an}} structure species of an
\xout{xx}\unsure{\hand{$\infty$}}-category  (besides the definition of standar\unsure{\over{t}{d}} fiber products)  -
without \unsure{\over{k}{h}}aving to get involved still less lost  in the technical intrica-%
cies of ever messier diagrams to write down,  with increasing order of
the structures to be added ...

\hand{6.}  In the outline of a method of construction for the structure species,
there has not been any explicit mention so far of the topological motiva-%
tion behind the whole approach,  which could wrongly give the impression
of being a purely algebraic one.  However,  topological considerations
alone are giving me the clue both for the description of the so-called
standar\unsure{\over{t}{d}} f\unsure{\over{u}{i}}ber products,  and of the inductive step allowing to wind up
from $\underline{C}_{n}$ to $\underline{C}_{n+1}$ ?  The heuristics \over{\unsure{though}}{indeed} of the present approach is sim-%
ple enough,  and suggested by the starting task,  to define pertinent
structure on the \inserted{system of} sets $F_{i}(X)$ of ``homotopies'' of arbitrary order\hand{,} associated
to an arbitrary topological space.  In effect,  the \inserted{covariant} functors  \xout{XxxxX}
\[ X \mapsto F_{i}(X) \ : \hand{(Spaces) $\longrightarrow$ (Sets)} \]
are representable by spaces $D_{i}$ ,  which are easily seen to be i-\deleted{\unsure{balls}}\inserted{cells}.
The source and  target maps  $s^{i}_{l}$ , $t^{i}_{l}$ : $F_{i}(X) \rightrightarrows F_{i-1}(X)$  are transposed
to maps, which I may denote by the same \unsure{letters}
\[ s^{i}_{l},\ t^{i}_{l} \ : \ D_{i-1} \rightrightarrows D_{i} \ . \]
Handling around a little,  one easily convinces oneself that all the
main structural items on $F_{*}(X)$ which one is figuring out in succession,
such as the degeneracy maps  $k^{i}_{l}$ ,  the inversion maps $inv_{i}$ ,  the composi
tion  $vou = v \underset{i}{*} u$  for  i-objects,  etc,  are all transposed of similar
maps which are defined between the cel\xout{\unsure{lex}}ls $D_{i}$ ,  or which go from such
cells to certain spaces,  deduced from these by gluing them together -
the most evident example \inserted[typed]{in this respect} being the composition of paths,  which is trans-%
posed of a map from the unit segment \hand{$I=$}$D_{1}$ into $(I,1) \amalg_{e} (I,0)$ ,  having pre-%
assigned values on the endpoints of I \hand{(} which correspond in fact to the
images of the two maps  $s^{1}_{1}$ , $t_{1}^{1}$ :  $D_{0}$ = one point $\rightrightarrows D_{1} = I$ \hand{)}.

\margin{24}
